\chapter*{TD 8 - Exercices Complet}

\setcounter{exercice}{0}

\begin{mdframed}
\begin{exercice}
Soit $X_1, \dots, X_n$ un échantillon i.i.d. de loi exponentielle de densité donnée pour tout $x \in \mathbb{R}$ par
\begin{equation*}
f_\lambda(x) = \lambda e^{-\lambda x} \mathbf{1}_{[0, +\infty[}(x),
\end{equation*}
de paramètre $\lambda > 0$ inconnu.

\begin{enumerate}
    \item Calculer la fonction de répartition de $X_1$.
    \item On note $m$ la médiane de $X_1$. Montrer que $m = \ln(2)/\lambda$.
    \item On note $\hat{m}_n$ la médiane empirique de l'échantillon. Rappeler la définition de $\hat{m}_n$ et donner la loi limite de $\sqrt{n}(\hat{m}_n - m)$ lorsque $n \to \infty$.
    \item On pose $\hat{\lambda}_n = \ln(2)/\hat{m}_n$. Donner la loi limite de $\sqrt{n}(\hat{\lambda}_n - \lambda)$.
\end{enumerate}
\end{exercice}
\end{mdframed}

\smallskip

\begin{mdframed}
\begin{exercice}
Soit $0 < a < b$ et $\theta \in [0, 1]$, on considère un $n$-échantillon $(X_1, \dots, X_n)$ de loi ayant pour densité $p_\theta$ définie pour tout $x \in \mathbb{R}$, par
\begin{equation*}
p_\theta(x) = \theta \frac{1}{a}\mathbf{1}_{[0, a]}(x) + (1 - \theta) \frac{1}{b}\mathbf{1}_{[0, b]}(x).
\end{equation*}
Dans la suite, on suppose que les paramètres $a$ et $b$ sont connus et on cherche à estimer le paramètre $\theta$ inconnu.

\begin{enumerate}
    \item Montrer que la vraisemblance est donnée par
    \begin{equation*}
    L(X_1, \dots, X_n; \theta) = \left( \frac{\theta}{a} + \frac{1 - \theta}{b} \right)^N \left( \frac{1 - \theta}{b} \right)^{n - N},
    \end{equation*}
    où $N = \text{Card}\{i \in \{1, \dots, n\} : X_i \in [0, a]\}$.
    \item Déterminer l'estimateur $\hat{\theta}$ du maximum de vraisemblance et calculer son biais.
    \item Calculer $\mathbb{E}[X_1]$ et en déduire un estimateur $\tilde{\theta}$ construit à l'aide de la méthode des moments et calculer son biais.
    \item Calculer le risque quadratique des estimateurs $\hat{\theta}$ et $\tilde{\theta}$. Comparer ces risques quadratiques quand $b = 2a$.
\end{enumerate}
\end{exercice}
\end{mdframed}

\smallskip

\begin{mdframed}
\begin{exercice}
Pour tous $a \in \mathbb{R}$ et $b > 0$, on note $\mathcal{P}(a, b)$ la loi sur $\mathbb{R}$ ayant pour densité
\begin{equation*}
p_{a,b}(x) = \frac{b}{(x - a)^{b+1}} \mathbf{1}_{[a+1, +\infty[}(x).
\end{equation*}

\subsection*{Partie I : Préliminaires}
\begin{enumerate}
    \item Vérifier que $p_{a,b}$ est bien une densité de probabilité sur $\mathbb{R}$.
    \item Calculer la fonction de répartition $F_{a,b}$ de $\mathcal{P}(a, b)$.
    \item Montrer que si $X \sim \mathcal{P}(a, b)$, alors $\ln(X - a)$ suit une loi exponentielle de paramètre $b$.
\end{enumerate}

\subsection*{Partie II : Estimation de $b$ dans le cas où $a$ est connu}
Dans cette partie, on considère un $n$-échantillon $(X_1, \dots, X_n)$ de loi $\mathcal{P}(a, b)$ et on suppose que $a$ est connu et $b$ inconnu.
\begin{enumerate}
    \setcounter{enumi}{3}
    \item Montrer que l'estimateur du maximum de vraisemblance de $b$ est donné par
    \begin{equation*}
    \hat{b}_n = \frac{n}{\sum_{i=1}^n \ln(X_i - a)}.
    \end{equation*}
    \item Justifier que $\hat{b}_n$ est consistant.
    \item Donner la limite en loi de $\sqrt{n}(\hat{b}_n - b)$ lorsque $n \to + \infty$.
\end{enumerate}

\subsection*{Partie III : Estimation de $a$ dans le cas où $b$ est connu}
Dans cette partie, on suppose que $a$ est inconnu et $b$ est connu et on cherche à estimer le paramètre $a$.
\begin{enumerate}
    \setcounter{enumi}{6}
    \item Montrer que l'estimateur du maximum de vraisemblance de $a$ est donné par
    \begin{equation*}
    \hat{a}_n = \min(X_1, \dots, X_n) - 1.
    \end{equation*}
    \item Calculer la fonction de répartition de $\hat{a}_n - a$. En déduire que $\hat{a}_n$ est consistant.
    \item En utilisant la question précédente, donner un intervalle de confiance non asymptotique de niveau $1 - \alpha$, $\alpha \in \, ]0, 1[$, pour estimer $a$.
\end{enumerate}
\end{exercice}
\end{mdframed}

\smallskip

\begin{mdframed}
\begin{exercice}
Soient $X_1, \dots, X_n$, des variables aléatoires i.i.d. de densité donnée pour tout $x \in \mathbb{R}$ par
\begin{equation*}
f(x) = \frac{\theta}{x^{\theta+1}} \mathbf{1}_{\{x \ge 1\}},
\end{equation*}
où $\theta$ est un paramètre inconnu.

\begin{enumerate}
    \item Dans cette question, on suppose d'abord que l'ensemble des paramètres est $\Theta = \{\theta > 1\}$. Calculer $\mathbb{E}_\theta[X_1]$ puis estimer $\theta$ par la méthode des moments.
    \item Dans cette question, on suppose désormais que l'ensemble des paramètres est $\Theta = \{\theta > 0\}$.
    \begin{enumerate}
        \item[i--] Montrer qu'il n'existe pas de $r > 0$ tel que la fonction $\theta \in \mathbb{R}_+^* \mapsto \mathbb{E}_\theta[X_1^r]$ soit injective.
        \item[ii--] Montrer que $\mathbb{E}_\theta\left[ \frac{1}{X_1} \right] = \frac{\theta}{\theta+1}$ puis estimer $\theta$ par la méthode des moments. On notera $\hat{\theta}_n^{MM}$ l'estimateur obtenu.
        \item[iii--] $\hat{\theta}_n^{MM}$ est-il consistant ?
        \item[iv--] Donner la loi asymptotique de $\sqrt{n}(\hat{\theta}_n^{MM} - \theta)$ lorsque $n \to +\infty$.
        \item[v--] Montrer que l'estimateur du maximum de vraisemblance $\hat{\theta}_n^{MV}$ vaut
        \begin{equation*}
        \hat{\theta}_n^{MV} = \frac{n}{\sum_{i=1}^n \ln(X_i)}.
        \end{equation*}
        \item[vi--] $\hat{\theta}_n^{MV}$ est-il consistant ?\\
        \textit{Indication : on pourra utiliser le changement de variables $u = \ln(x)$.}
    \end{enumerate}
\end{enumerate}
\end{exercice}
\end{mdframed}

\smallskip

\begin{mdframed}
\begin{exercice}
Soient $X_1, \dots, X_n$ des variables aléatoires i.i.d. de densité donnée pour tout $x \in \mathbb{R}$ par
\begin{equation*}
f(x) = \frac{1}{2\sqrt{x\theta}} \mathbf{1}_{]0, \theta]}(x),
\end{equation*}
où $\theta > 0$ est un paramètre inconnu.

\begin{enumerate}
    \item Calculer $\mathbb{E}_\theta[X_1]$.
    \item En déduire un estimateur $\hat{\theta}_n$ par la méthode des moments.
    \item Déterminer l'estimateur $\tilde{\theta}_n$ de $\theta$ par maximum de vraisemblance.
    \item Déterminer la loi de $\tilde{\theta}_n$.
    \item Montrer que le risque quadratique de $\tilde{\theta}_n$ tend vers 0 quand $n$ tend vers $+\infty$.
    \item Comparer les risques quadratiques de $\hat{\theta}_n$ et de $\tilde{\theta}_n$.
    \item Préciser la loi asymptotique de $\sqrt{n}(\hat{\theta}_n - \theta)$ et de $n(\theta - \tilde{\theta}_n)$.
    \item Montrer que la médiane de $X_1$ est $F_\theta^{-1}(1/2) = \theta/4$. En déduire un estimateur $\hat{\theta}_n^m$ consistant pour l'estimation de $\theta$.
    \item Donner la loi limite de $\hat{\theta}_n^m - \theta$ après une renormalisation convenable.
\end{enumerate}
\end{exercice}
\end{mdframed}

\smallskip

\begin{mdframed}
\begin{exercice}
Soient $X_1, \dots, X_n$ des variables aléatoires i.i.d. de densité donnée pour tout $x \in \mathbb{R}$ par
\begin{equation*}
f(x) = e^{-\theta} \mathbf{1}_{]0, e^\theta]}(x),
\end{equation*}
où $\theta > 0$ est un paramètre inconnu.

\begin{enumerate}
    \item Déterminer l'estimateur du maximum de vraisemblance de $\theta$.
    \item Montrer que $\theta - \ln(X_{(n)})$ suit une loi exponentielle de paramètre $n$ (où $X_{(n)} = \max(X_1, \dots, X_n)$).
    \item Soit $0 < \alpha < 1$. Montrer que 
    \begin{equation*}
    \left[ \ln(X_{(n)}) \, , \, \ln(X_{(n)}) - \frac{1}{n} \ln(\alpha) \right]
    \end{equation*}
    est un intervalle de confiance pour $\theta$ de niveau $1 - \alpha$.
\end{enumerate}
\end{exercice}
\end{mdframed}

\smallskip

\begin{mdframed}
\begin{exercice}

\textbf{Rappels}

\smallskip

\begin{itemize}
\item Une v.a.  de loi exponentielle $A \sim \mathcal{E}(\theta)$ , $\theta >0$ admet pour densité :

$$
f_A(x) = \theta e^{-\theta x} \mathbb{1}_{x>0}
$$
\item Soit $A_1,\cdots,A_n$ un $n$-échantillon i.i.d. de loi $\mathcal{E}(\theta)$. Alors $S = \sum_{i=1}^n A_i$ suit une loi Gamma $\Gamma(n,\theta)$ de densité donnée par

$$
f_S(x) = \frac{\theta^n}{\Gamma(n)} x^{n-1} e^{- \theta x} \mathbb{1}_{x>0}
$$

où $\Gamma$ désigne la fonction Gamma. 
\item Pour tout entier entier $n \in \mathbb{N}^*$, $\Gamma(n) = (n-1)!$
\end{itemize}

\smallskip

\textbf{Questions préliminaires}

\smallskip

\begin{enumerate}
\item Soit $A \sim \mathcal{E}(\theta)$ avec $\theta > 0$. Soit $A_1, \cdots A_n$, $n$ copies i.i.d. de $A$. Donner la vraisemblance associée aux observation $A_1, \cdots A_n$.
\item Soit $B \sim \mathcal{B}(p)$ avec $p \in ]0,1[$. Soit $B_1, \cdots B_n$, $n$ copies i.i.d. de $B$. Donner la vraisemblance associée aux observation $B_1, \cdots B_n$.
\end{enumerate}

\smallskip

\textbf{Problème}

\smallskip

On s'interesse à la durée de vie de deux composants électroniques se trouvant dans un système solidaire. Si l'un des deux composants tombe en panne, le système tout entier doit être changé.

Les durée de vie des deux composants sont modélisés par des variables aléatoires exponentielles de paramètres $\lambda > 0$ et $\mu > 0$ indépendantes. Pour $n$ systèmes, on a donc un $n$-échatillon i.i.d. $X_1,\cdots,X_n$ de loi $\mathcal{E}(\lambda)$ et un $n$-échantillon i.i.d. $Y_1, \cdots , Y_n$ de loi $\mathcal{E}(\mu)$, indépendant du précédent.

Cépendant, on observe uniquement la durée de vie du commposant tombé en panne, ce qui veut dire que l'on observe uniquement le $n$-échantillon 

$(Z_1,W_1) , \cdots , (Z_n,W_n)$, où pour tout $i \in 1 , \cdots , n$,

$$
Z_i = \min (X_i,Y_i) \qquad W_i = \mathbb{1}_{Z_i = X_i} = \mathbb{1}_{X_i \leq Y_i}   
$$

\smallskip

\begin{enumerate}
\item Calculer la fonction de répartition de $Z_i$ et en déduire $Z_i \sim \mathcal{E} (\lambda + \mu )$
\item Montrer que $W_i \sim \mathcal{B}( \frac{\lambda}{\lambda + \mu} )$
\item
\begin{enumerate}
\item Soit $z \geq 0$. Montrer que $\mathbb{P}(Z_i \leq z , W_i = 1) = \frac{\lambda}{\lambda + \mu} \mathbb{P}(Z_i \leq z_i)$.
\item En conclure que les v.a. $Z_i$ et $W_i$ sont indépendantes.
\end{enumerate}
\end{enumerate}

On rappelle que par indépendance entre les $Z_i$ et les $X_i$, la vraisemblance du $n$-échantillon i.i.d. $(Z_1,W_1) , \cdots , (Z_n,W_n)$ est définie comme le produit de la vraisemblance du $n$-échantillon $Z_1 , \cdots , Z_n$ par la vraisemblance du $n$-échantillon $W_1 , \cdots , W_n$. On suppose désormais que le paramètre $\mu$ \textbf{est connu.}

\begin{enumerate}
\item Déterminer l'EMV $\hat{\lambda}_n$ de $\lambda$.
\item Montrer que $\hat{\lambda}_n$ est consistant.
\item Calculer le biais de $\hat{\lambda}_n$ puis montrer qu'il tend ver 0 lorsque $n \to \infty$.
\item Déterminer la loi asymptotique de $\sqrt{n} ( \hat{\lambda}_n - \lambda)$ lorsque $n \to \infty$.
\item Donner un intervalle de confiance bilatéral symétrique asymptotique de niveau $1 - \alpha$ pour $\lambda$. Dans la question suivantee on suppose que $\lambda$ et $\mu$ \textbf{sont inconnus.}
\item Déterminer l'EMV $(\hat{\lambda}_n,\hat{\mu}_n)$ de $( \lambda , \mu)$.
\end{enumerate}
\end{exercice}
\end{mdframed}

\smallskip

\begin{mdframed}
\begin{exercice}
Soit $X_1, \cdots, X_n$ des variables aléatoires i.i.d. de deinsité donnée, pour tout $x \in \mathbb{R}$, par :

$$
f_\theta(x) = C_\theta x^2 \mathbb{1}_{]0,\theta]} (x)
$$

où $\theta > 0$ est un paramètre inconnu, et $C_\theta$ une constante ne dépendant que de $\theta$.

\begin{enumerate}
\item Justifier que $C_\theta = \frac{3}{\theta^3}$.
\item 
\begin{enumerate}
\item Calculer $\mathbb{E}_\theta [X_1]$
\item En déduire un estimateur $\hat{\theta}^{MM}_n$ par la méthode des moments.
\item Calculer le risque quadratique de $\hat{\theta}^{MM}_n$.
\item Donner la loi limite de $\hat{\theta}^{MM}_n - \theta$ après une renormalisation convenable.
\end{enumerate}
\item Calculer la fonction de répartition de $X_1$ notée $F_\theta$.
\item
\begin{enumerate}
\item Montrer que la médiane de $X_1$ est $x_{\frac{1}{2}} := F_\theta^{-1} (\frac{1}{2}) = 2^{-\frac{1}{3}} \theta$. En déduire une estimateur $\hat{\theta}^{m}_n$ consistant pour l'estimation de $\theta$.
\item Donner la loi limite de $\hat{\theta}^{m}_n - \theta$ après une renormalisation convenable.
\end{enumerate}
\item
\begin{enumerate}
\item Déterminer l'estimateur $\hat{\theta}^{MV}_n$ de $\theta$ par maximum de vraisemblance.
\item Déterminer la fonction de répartition de $v\hat{\theta}^{MV}_n$.
\item Montrer que $n(\theta-\hat{\theta}^{MV}_n)$ tend en loi vers une loi exponentielle dont on précisera le paramètre.
\end{enumerate}
\item Parmi les trois estimateurs, lequel faut-il privilégeir et porquoi?
\item Soit $\alpha \in ]0,1[$ Utiliser l'estimateur choisi pour construire un intervalle de confiance (asymptotique ou non) de niveau de condiance $1-\alpha$.
\end{enumerate}
\end{exercice}
\end{mdframed}